\section{Classical Lambek calculus and completeness}
The classical literature provides an important baseline for this thesis. The original system of
\citet{Lambek1958} and \citet{Lambek1961} was quickly recognised as more than a linguistic
notation: it is a mathematically robust substructural logic with a rich proof theory and model
theory. Surveys and handbook treatments such as \citet{vanBenthem1988} and
\citet{Moortgat1997} show how the calculus sits within the wider landscape of categorial and
type-logical grammar.

For the specific question of completeness, the paper by \citet{Dosen1985} is especially relevant
because it directly establishes a completeness theorem for the Lambek calculus of syntactic
categories. Equally important is the relational-semantics perspective of
\citet{AndrekaMikulas1994}, who showed that relational models yield strong completeness results
for suitable versions of the calculus. From a language-theoretic viewpoint, \citet{Pentus1995}
proved completeness with respect to language models, while \citet{Pentus1993} and
\citet{Pentus1997} clarified the relation between Lambek grammars and context-free grammars.

These results matter for the present thesis in two ways. First, they demonstrate that
completeness questions for Lambek systems are deep but tractable. Second, they show that
there are already several mathematically natural semantic classes for which completeness holds.
This makes it reasonable to ask whether a similar story can be told for a probabilistic,
corpus-based semantics.

